\subsection{Universal robots' robots}
\subsubsection{Interface used by SMC}
We use the \href{https://sdurobotics.gitlab.io/ur_rtde/introduction/introduction.html}{ur\_rtde}
implementation of the \href{https://www.universal-robots.com/developer/communication-protocol/rtde/}{RTDE protocol}.
It is very straightforward and works exactly as expected for SMC purposes and API.
The particulars can be found in table \ref{table-smc-ur-rtde-interface}.
At a higher level, the robot opens a TCP socket dedicated to the RTDE protocol 
at port 30002(? could be +/- 1 or 2, doesn't matter really).
ur\_rtde creates a client side TCP socket which connects to the robot's RTDE socket.
All socket calls, message serialization etc is handled by ur\_rtde and it works without faults.
Importantly,  ur\_rtde fetches all data from the robot at every time-step,
and stores it in its internal struct. This means that there is no communication
cost to accessing particular data.

On the usage side, despite there being only 1 client side socket, there are 3 classes for using it:
RTDEControlInterface, RTDEReceiveInterface and RTDEIOInterface.
They all share the same socket, and this is just for better code design
(separation of concerns, obvious from the name).

\begin{table}[h!]
		\centering
		\label{table-smc-ur-rtde-interface}
		\begin{tabular}{p{4.5cm} | p{8cm}}
				\textbf{SMC function} & \textbf{How it's achieved with ur\_rtde} \\
				\hline
				connectToRobot() & \begin{itemize}
						\item RTDEControlInterface(args.robot\_ip)
						\item RTDEReceiveInterface(args.robot\_ip)
						\item RTDEIOInterface(args.robot\_ip)
				\end{itemize} \\
				\hline
\_updateQ() & rtde\_receive.getActualQ() \\
				\hline
\_updateV() & rtde\_receive.getActualQd() \\ 
				\hline
sendVelocityCommandToReal() & rtde\_control.speedJ(v, args.acceleration, 1/args.ctrl\_freq)
\newline
\textbf{IMPORTANT:} acceleration is a scalar, same for all joints (effectively).
dt is for how long the robot will try to achieve the prescribe velocity v,
i.e. it's the hangup time for the command. 
We set it to be the specified control frequency (500Hz max).\\
				\hline
\_updateWrench() & rtde\_receive.getActualTCPForce(). 
\newline
\textbf{IMPORTANT:} This returns the unfiltered wrench in the \textbf{base frame} of 
the robot. We immediately map it to the end-effector frame $ T_{ w }^{ e }  $.
Upon the start of every script, we also first collect samples to collect the bias,
and offset it. Despite calling rtde\_control.zeroFtSensor(), there can be bias.
And it builds up quickly. After a few minutes it should be recalibrated
if running a continuous experiment.
The sensor is noisy, our baseline filtering gets reliable wrench readings
if the norm of the wrench is \textbf{at least 2N}.\\
setInitialPose() & rtde\_receive.getActualQ() if connected to the robot
and with --no-real (used to match sim to reality before experiment)\\
				\hline
stopRobot() & first send zeros velocity, then go to freedrive. 
Without going to freedrive, the robot keeps trying to go at zero velocity,
but in actuality keeps moving slowly (collides with something in 5-10min)\\
				\hline
connectToGripper() & gripper dependent, look below if connected with flange to arm's wrist.\\
				\hline
setFreedrive() & rtde\_control.freedriveMode() \\
				\hline
unSetFreedrive() & rtde\_control.endFreedriveMode() \\
				\hline
\_sendTorqueCommand() & is available on UR robots now, but there are
no torque-level controllers in SMC. This is future work.
		\end{tabular}
		\caption{SMC-ur\_rtde interface.}
\end{table}

\subsubsection{An opinionated remark on alternative interfaces}
This remark is a ``heads up'' for those who for whatever reason
have to read UR official documentation.

UR has developed its own programming language for controlling
their robots, called URscript. It is simply a butchered version of Python.
As in, it is literally Python, but heavily constricted (can't install Python packages to it for example),
with premade functions to control the robot.
There are 2 additional sockets to control the robot with remotely.
One accepts entire URscript programs and then executes them.
The other waits for one URscript command at a time, and executes it on completion.

Both of these options are completely useless for the purpose of experimental control engineering
because URscript is not a real programming language.
They are very useful for technicians to safely used the robot with the functionality
developed by UR. This is not SMC's target audience. The target audience are people
who want to test novel controllers.
Avoid non-RTDE interfaces completely. Use available SMC's compliance control, 
or even better, write your own (it will cost you less time and get a better result
than reverse-engineering a URscript program you found.)


\subsubsection{Grippers supported in SMC}
\begin{itemize}
		\item some Robotiq grippers (Hand-E, 2FG), via flange
		\item OnRobot  TwoFg (external cabling)
		\item (soon) Schunk EGK-40 via flange
\end{itemize}

\subsubsection{Notes on integrating grippers connected via the flange}
There are 2 distinct, mutually exclusive options:
\begin{enumerate}
		\item communicate with the gripper directly, with the gripper protocol, over the TCP socket
exposed by the robot
		\item find a URcap for your gripper, and find Python code to interface
			with whatever the URcap expects
\end{enumerate}
\paragraph{Direct communication with the gripper}
The robot opens a TCP socket on port 63352 for direct communication with the gripper.
My understanding is that whatever bytes you send to that socket
will be exactly relayed to the gripper.
There is nothing special about the client side socket you construct.
Typical industrial gripper are very simple from a control standpoint,
and so is communication with them. Namely, the state of the gripper
is likely less than 20 bytes in size. The commands you give it are a few bytes
in length (ex. 1 bytes for target position, 1 for motion velocity, and 1 for grasping force).
My recommendation is to read the gripper's documentation, which has to include
which bytes mean what, and implement the communication protocol directly,
with raw bytes being sent and read. It is very likely you will need to read the gripper's
documentation if you just bought it (preferably you read it \textit{before} you bought it).
(Limited) Experience shows that this is the most stable, set it and forget it way to control the gripper.

\paragraph{URcap solution}
If whoever provides the URcap also provides Python code to interact with it,
this is the best possible solution.
However, chances are that there is no Python code. Be wary if you are determined
to find a ``simple'' workaround. This is a dark path to tread on.
I have seen grippers which provide web interfaces, and people writing Python code
to trigger Javascript to move the GUI sliders.
I suggest toughing it out and writing the raw bytes yourself, that way there is 
at least a clear end in sight. But there is no winning here, only tough labour, or
implementing hacks you wish you can hide and forget, or even worse --- external cabling.

